\subsection{Documentation}
Software documentation \cite{naimi2024automate_4} is not just a tool, but a crucial link that connects developers, end users, and maintainers. It serves as a beacon, providing essential information on the software systems, guiding their creation, understanding, and maintenance \cite{naimi2024automate_5}. Internal documentation \cite{naimi2024automate_6} is more than just a communication tool in the software development process. It's a lifeline, helping to maintain clear guidelines and responsibilities, fostering efficient collaboration among the development teams.Comprehensive documentation is a game-changer for external stakeholders, such as user experience professionals and end users. User documentation is not just a guide, but a key to understanding the software's functionality and troubleshooting procedures. It empowers users to navigate a developed software product, enhancing overall user satisfaction \cite{naimi2024automate_7}.

In their paper about creating documentation from UML diagrams, \citet{naimi2024automate} found that when compared to conventional methods like MDA \cite{naimi2024automate_mda}, this approach offers significant advantages. First, the quality of the generated documentation is remarkably human-like, providing clear and relevant explanations that closely resemble expertly written manuals. Secondly, they speed up the documentation process, which is very much needed by fast software development cycles.

\paragraph{Limitations}
The use of LLMs is generally limited to a certain number of tokens that can be used for free. Limited tokens can pose a challenge for projects with extensive documentation requirements. Furthermore, those models need a lot of fine-grained context to generate proper documentation. The more such context is provided, the better the output generally will be. That requires understanding the domain and effective communication with the language model \cite{naimi2024automate}. \citet{diggs2024leveragingllmslegacycode} states that LLMs still struggle with low-level or highly specialized legacy systems such as that of IBM Assembly Language Code (ALC).